% GENERATED FILE. Edit canonical lessons, then run npm run generate:books.
% canonical-source-hash: fnv1a64:2369a82285d724d3
% canonical-lessons: BN-C403-ghor, BN-C403-buking, BN-C403-lipht, BN-C403-risepshon, BN-C403-bondobosto-kora

\chapter{Checking In}
\label{ch:bn-checking-in}
\hlchaptermodality{403}

\begin{chapteropening}
\textbf{By the end of this chapter:} \emph{I can check in, collect a key and ask for something to be arranged.}

Complete the last lesson of chapter 403: I can check in, collect a key and ask for something to be arranged.
\end{chapteropening}

\section[\texorpdfstring{\bn{ঘর} (ghôr)}{ঘর (ghôr)}]{\bn{ঘর} (ghôr) — a room (in a house); a home}

\label{lesson:BN-C403-ghor}



\begin{quote}
Before the new one: say the Bengali for at a loss, then the Bengali for to go round in circles.
\end{quote}

\subsection*{You'll want to know: \bn{ঘর}}
\textbf{\bn{ঘর}} — \emph{ghôr} — \textquotedblleft{}a room (in a house); a home\textquotedblright{}.

\begin{quote}
\bn{হোটেলে} \bn{আমাদের} \bn{ঘরটা} \bn{ছোট}, \bn{কিন্তু} \bn{পরিচ্ছন্ন।}
\end{quote}

\begin{quote}
\emph{hoṭele āmāder ghôrṭā chhoṭo, kintu pôrichchhônnô.}
\end{quote}

\begin{quote}
Our room at the hotel is small but clean.
\end{quote}

\bn{ঘর} is any room, and home as well: \bn{ঘরে} \bn{ফেরা}, to come home. \bn{রুম} is the hotel word.

\begin{cousinweb}
\textbf{\bn{ঘর}} (\emph{ghôr}) is inherited: CDIAL lists Bengali \emph{ghar}, \textquotedblleft{}a house\textquotedblright{}, under \emph{ghara}, \textquotedblleft{}a house\textquotedblright{} (CDIAL 4428), and notes that the usual link with Sanskrit \emph{gṛha} is hard to make work by sound.
\end{cousinweb}

\begin{grammarlens}[title={What you've built}]
One more word for getting around.
\end{grammarlens}

\subsection*{Your turn}
\begin{itemize}
\raggedright
  \item \emph{Say it:} \emph{ghôr}
  \item \emph{Say it:} \emph{ghôr}, once more
  \item \emph{Say it:} \emph{hoṭele āmāder ghôrṭā chhoṭo, kintu pôrichchhônnô.} — Our room at the hotel is small but clean.
  \item \emph{Recall:} say the Bengali for at a loss, then the Bengali for to go round in circles, then say \emph{ghôr} again
\end{itemize}

\begin{tcolorbox}[breakable,colback=teal!4,colframe=teal!35!black,title={Before you move on}]
What does \bn{ঘর} mean? (\textquotedblleft{}A room (in a house); a home\textquotedblright{}.) Say it once more.
\end{tcolorbox}
\section[\texorpdfstring{\bn{বুকিং} (buking)}{বুকিং (buking)}]{\bn{বুকিং} (buking) — a booking; a reservation}

\label{lesson:BN-C403-buking}



\begin{quote}
Before the new one: say the Bengali for to go round in circles, then the Bengali for a room (in a house).
\end{quote}

\subsection*{You'll want to know: \bn{বুকিং}}
\textbf{\bn{বুকিং}} — \emph{buking} — \textquotedblleft{}a booking; a reservation\textquotedblright{}.

\begin{quote}
\bn{আমরা} \bn{এক} \bn{সপ্তাহ} \bn{আগে} \bn{হোটেলের} \bn{বুকিং} \bn{করলাম।}
\end{quote}

\begin{quote}
\emph{āmrā ek sôptāhô āge hoṭeler buking kôrlām.}
\end{quote}

\begin{quote}
We made the hotel booking a week ahead.
\end{quote}

\begin{cousinweb}
\textbf{\bn{বুকিং}} (\emph{buking}) is borrowed from English \emph{booking}; \bn{বুক} \bn{করা}, to book, is heard too, but \bn{বুক} alone is the chest.
\end{cousinweb}

\begin{grammarlens}[title={What you've built}]
One more word for getting around.
\end{grammarlens}

\subsection*{Your turn}
\begin{itemize}
\raggedright
  \item \emph{Say it:} \emph{buking}
  \item \emph{Say it:} \emph{buking}, once more
  \item \emph{Say it:} \emph{āmrā ek sôptāhô āge hoṭeler buking kôrlām.} — We made the hotel booking a week ahead.
  \item \emph{Recall:} say the Bengali for to go round in circles, then the Bengali for a room (in a house), then say \emph{buking} again
\end{itemize}

\begin{tcolorbox}[breakable,colback=teal!4,colframe=teal!35!black,title={Before you move on}]
What does \bn{বুকিং} mean? (\textquotedblleft{}A booking; a reservation\textquotedblright{}.) Say it once more.
\end{tcolorbox}
\section[\texorpdfstring{\bn{লিফট} (liphṭ)}{লিফট (liphṭ)}]{\bn{লিফট} (liphṭ) — a lift}

\label{lesson:BN-C403-lipht}



\begin{quote}
Before the new one: say the Bengali for a room (in a house), then the Bengali for a booking.
\end{quote}

\subsection*{You'll want to know: \bn{লিফট}}
\textbf{\bn{লিফট}} — \emph{liphṭ} — \textquotedblleft{}a lift\textquotedblright{}.

\begin{quote}
\bn{লিফটটা} \bn{বন্ধ} \bn{ছিল}, \bn{তাই} \bn{আমরা} \bn{হেঁটে} \bn{উঠলাম।}
\end{quote}

\begin{quote}
\emph{liphṭṭā bôndhô chhilô, tāi āmrā hẽṭe uṭhlām.}
\end{quote}

\begin{quote}
The lift was out of order, so we walked up.
\end{quote}

\begin{cousinweb}
\textbf{\bn{লিফট}} (\emph{liphṭ}) is borrowed from English \emph{lift}, with \bn{ফ} for f and \bn{ট} for t.
\end{cousinweb}

\begin{grammarlens}[title={What you've built}]
One more word for getting around.
\end{grammarlens}

\subsection*{Your turn}
\begin{itemize}
\raggedright
  \item \emph{Say it:} \emph{liphṭ}
  \item \emph{Say it:} \emph{liphṭ}, once more
  \item \emph{Say it:} \emph{liphṭṭā bôndhô chhilô, tāi āmrā hẽṭe uṭhlām.} — The lift was out of order, so we walked up.
  \item \emph{Recall:} say the Bengali for a room (in a house), then the Bengali for a booking, then say \emph{liphṭ} again
\end{itemize}

\begin{tcolorbox}[breakable,colback=teal!4,colframe=teal!35!black,title={Before you move on}]
What does \bn{লিফট} mean? (\textquotedblleft{}A lift\textquotedblright{}.) Say it once more.
\end{tcolorbox}
\section[\texorpdfstring{\bn{রিসেপশন} (risepshôn)}{রিসেপশন (risepshôn)}]{\bn{রিসেপশন} (risepshôn) — a reception desk}

\label{lesson:BN-C403-risepshon}



\begin{quote}
Before the new one: say the Bengali for a booking, then the Bengali for a lift.
\end{quote}

\subsection*{You'll want to know: \bn{রিসেপশন}}
\textbf{\bn{রিসেপশন}} — \emph{risepshôn} — \textquotedblleft{}a reception desk\textquotedblright{}.

\begin{quote}
\bn{রিসেপশনে} \bn{আমরা} \bn{ঘরের} \bn{চাবি} \bn{নিলাম।}
\end{quote}

\begin{quote}
\emph{risepshône āmrā ghôrer chābi nilām.}
\end{quote}

\begin{quote}
We collected the room key at reception.
\end{quote}

\begin{cousinweb}
\textbf{\bn{রিসেপশন}} (\emph{risepshôn}) is borrowed from English \emph{reception}, the desk where guests arrive.
\end{cousinweb}

\begin{grammarlens}[title={What you've built}]
One more word for getting around.
\end{grammarlens}

\subsection*{Your turn}
\begin{itemize}
\raggedright
  \item \emph{Say it:} \emph{risepshôn}
  \item \emph{Say it:} \emph{risepshôn}, once more
  \item \emph{Say it:} \emph{risepshône āmrā ghôrer chābi nilām.} — We collected the room key at reception.
  \item \emph{Recall:} say the Bengali for a booking, then the Bengali for a lift, then say \emph{risepshôn} again
\end{itemize}

\begin{tcolorbox}[breakable,colback=teal!4,colframe=teal!35!black,title={Before you move on}]
What does \bn{রিসেপশন} mean? (\textquotedblleft{}A reception desk\textquotedblright{}.) Say it once more.
\end{tcolorbox}
\section[\texorpdfstring{\bn{বন্দোবস্ত} \bn{করা} (bôndobôstô kôrā)}{বন্দোবস্ত করা (bôndobôstô kôrā)}]{\bn{বন্দোবস্ত} \bn{করা} (bôndobôstô kôrā) — to arrange; to make arrangements}

\label{lesson:BN-C403-bondobosto-kora}



\begin{quote}
Before the new one: say the Bengali for a lift, then the Bengali for a reception desk.
\end{quote}

\subsection*{You'll want to know: \bn{বন্দোবস্ত} \bn{করা}}
\textbf{\bn{বন্দোবস্ত} \bn{করা}} — \emph{bôndobôstô kôrā} — \textquotedblleft{}to arrange; to make arrangements\textquotedblright{}.

\begin{quote}
\bn{মামা} \bn{আমাদের} \bn{থাকার} \bn{বন্দোবস্ত} \bn{করলেন।}
\end{quote}

\begin{quote}
\emph{māmā āmāder thākār bôndobôstô kôrlen.}
\end{quote}

\begin{quote}
My uncle arranged somewhere for us to stay.
\end{quote}

Past (he/she …): \bn{তিনি} \bn{গাড়ির} \bn{বন্দোবস্ত} \bn{করলেন}, \emph{tini gāṛir bôndobôstô kôrlen}.

The noun alone, \bn{বন্দোবস্ত}, is an arrangement; \bn{ব্যবস্থা} says much the same in a more careful way.

\begin{cousinweb}
\textbf{\bn{বন্দোবস্ত} \bn{করা}} (\emph{bôndobôstô kôrā}) joins \bn{বন্দোবস্ত}, usually traced to Persian \emph{band-o-bast}, \textquotedblleft{}tying and binding\textquotedblright{}, a settlement of affairs, to \bn{করা}, to do.
\end{cousinweb}

\begin{grammarlens}[title={What you've built}]
One more word for getting around.
\end{grammarlens}

\subsection*{Your turn}
\begin{itemize}
\raggedright
  \item \emph{Say it:} \emph{bôndobôstô kôrā}
  \item \emph{Say it:} \emph{bôndobôstô kôrā}, once more
  \item \emph{Say it:} \emph{māmā āmāder thākār bôndobôstô kôrlen.} — My uncle arranged somewhere for us to stay.
  \item \emph{Recall:} say the Bengali for a lift, then the Bengali for a reception desk, then say \emph{bôndobôstô kôrā} again
\end{itemize}

\begin{tcolorbox}[breakable,colback=teal!4,colframe=teal!35!black,title={Before you move on}]
What does \bn{বন্দোবস্ত} \bn{করা} mean? (\textquotedblleft{}To arrange; to make arrangements\textquotedblright{}.) Say it once more.
\end{tcolorbox}
